\chapter{Introduction}
\label{ch:introduction}

During a multi-stage incident, an analyst may need to decide which
activity to investigate before the next attacker action is visible. A
forecast can prioritize this investigation by assigning probabilities
to plausible follow-up actions. It is important that the forecast
reflect the evidence that could actually have reached the analyst or
detector. Audit records carry an occurrence time, while their availability
to a predictor depends on when the collector receives them.

Consider a sequence of file-related activities involving two targets.
A termination record for one process may arrive after another process
starts. Benign activity on the host may add further start and termination
records in between. A model that treats the most recent exit code as the
result relevant to every subsequent decision can confuse the two
targets. A model that reconstructs a past prefix using all audit
records by occurrence time can also inadvertently give the predictor
information that the collector had yet to receive.

This thesis studies \textbf{next actual attacker action prediction}
from Windows audit logs. The label is the semantic category of the
first source-designated attacker action that actually starts after
the prediction boundary. The output is a distribution over 18 action
categories. An execution can provide several prediction rows, each
with its own evidence prefix and label. The task is evaluated at known
decision boundaries in a controlled environment. Inferring those
boundaries in an operational deployment is a separate problem.

PROPHET represents the evidence at each boundary in two complementary
forms. A gated recurrent unit summarizes the events in receipt order.
A target adapter connects observed events through process identity and
observed target membership. These links allow the predictor to use
the result of an operation on the relevant target even when unrelated
events intervene. A target-free path produces a prediction when the
observed prefix supplies no usable target relation.

The main question is whether linking events to the correct target
improves the forecast. On 64 new runs of previously evaluated
configurations, relation attention reduced execution-weighted negative
log likelihood (NLL) from 0.2171 to 0.0458 compared with a sequence-only
GRU. In a separate comparison, the target history model used the same
events and the same group sizes while event membership was exchanged.
On 127 previously examined executions, this exchange increased NLL
from 0.0584 to 0.0869. The first result measures the gain of the
target-aware predictor as a whole. The second tests whether correct
membership contributes to that gain in the target history model.

Chapter~\ref{ch:related} places these comparisons in the context of ATT\&CK, attack graphs, and prior forecasting systems.

The thesis defines an evaluation in which every forecast uses records
already received at its decision boundary. It tests observed target
membership against sequence and grouping controls, including new
executions with unfamiliar target names. It then tests how to present
recurrent history to the readout when benign activity is interleaved.
Chapter~\ref{ch:task} follows a recorded prediction to show what
the target links carry. Chapter~\ref{ch:method} defines the models,
and Chapters~\ref{ch:evaluation} and~\ref{ch:results} report the
comparisons and the conditions under which the gains hold.
